\section*{Mittelalterliche Rahmfladen}
\ihead{}\chead{Personen:1}\ohead{}
\ifoot{Vorbereitungszeit:20 min}\cfoot{}\ofoot{Kochzeit:30 min}
{\Large Zutaten}
\begin{multicols}{2}
\begin{itemize}
\item \SI{10}{g} Hefe
\item etwas lauwarmes Wasser
\item etwas Mehl
\item \SI{150}{g} Mehl
\item \SI{70}{g} Roggenmehl
\item \SI{1,5}{TL} Salz
\item \SI{150}{ml} lauwarmes Wasser
\item Fett für das Blech
\item Crème fraîche, ca. \SI{2}{EL} pro Fladen
\item Speckwürfel nach Bedarf
\item Lauchzwiebeln, in Ringe geschnitten, nach Bedarf
\end{itemize}
% \columnbreak
\end{multicols}
\noindent {\Large Zubereitung}\\
\\
\textit{Vorbereitung:} Zuerst mit der Hefe, etwas Wasser und etwas Mehl einen Vorteig ansetzen.
Dazu die Hefe in ein kleines Schälchen bröseln und mit ca. \SI{2}{EL} Mehl und warmem Wasser zu einem glatten Teig verrühren.
Den Vorteig stehen lassen, bis er Bläschen schlägt.
In der Zwischenzeit die beiden Mehlsorten mit dem Salz und dem Wasser mischen.
Den Vorteig zur Mehlmischung geben und alles gut verkneten.
Der Teig ist fertig, wenn er nicht mehr an den Händen klebt und sich gut als Klumpen aus der Schüssel lösen lässt.
Sollte der Teig zu lange kleben, etwas Mehl hinzufügen.
Sollten Mehlreste in der Schüssel bleiben, etwas Wasser hinzufügen.
Die Schüssel mit einem Handtuch abdecken und den Teig ca. \SI{1}{h} ruhen lassen.
Anschließend den Teig zu Fladen formen und auf ein eingefettetes Blech legen.
Die Fladen mit Crème fraîche bestreichen und mit Speckwürfeln bestreuen.
Die Fladen im vorgeheizten Backofen bei ca. \SI{200}{\celsius} Ober-/Unterhitze \SI{25-30}{min} backen.
Die fertig gebackenen Fladen anschließend mit Lauchzwiebelringen bzw. Zwiebelringen bestreuen.
Nach Belieben etwas Pfeffer darüberstreuen.
